\section*{Capítulo \ref{ch:b2:anthropogenic-climate-change} --- El cambio climático antropógeno}

\begin{solution}{exo:b2:anthropogenic-climate-change:1}
(a) $3.75 - 0.60 + 0.22 + 0.45 + 0.11 = \qty{3.93}{W.m^{-2}}$, ya que los ajustes suman
\qty{0.18}{W.m^{-2}}. (b) $3.93/3.71 = 1.059$: un \qty{5.9}{\percent} más. (c) El ajuste de la temperatura
troposférica, $-0.60$: una troposfera más caliente emite más infrarrojo hacia el espacio, lo que retira
parte de la energía adicional antes de que la superficie haya respondido.
\end{solution}

\begin{solution}{exo:b2:anthropogenic-climate-change:2}
(a) $4\pi(\num{6.371e6})^2 = \qty{5.10e14}{m^2}$. (b) En 47 años, $\num{434.9e21}/(47 \times
\num{3.156e7} \times \num{5.10e14}) = \qty{0.575}{W.m^{-2}}$; en 12 años, $\num{152.4e21}/(12
\times \num{3.156e7} \times \num{5.10e14}) = \qty{0.789}{W.m^{-2}}$, los 0.57 y 0.79 de la
evaluación. (c) $0.91 \times 434.9 = \qty{396}{ZJ}$.
\end{solution}

\begin{solution}{exo:b2:anthropogenic-climate-change:3}
(a) $-3.22 + 1.30 + 0.35 + 0.42 - 0.01 = \num{-1.16}$, la neta. (b) $3.93/1.16 =
\qty{3.39}{\celsius}$; $3.93/3.22 = \qty{1.22}{\celsius}$. (c) $3.39/1.22 = 2.78$: las retroalimentaciones
casi triplican el calentamiento de Planck.
\end{solution}

\begin{solution}{exo:b2:anthropogenic-climate-change:4}
(a) Un petagramo de carbono son $44.009/12.011 = 3.664$~Gt de dióxido de carbono: $0.45 \times 3.664 =
\qty{1.65}{\celsius}$ por \qty{1000}{PgC}, y el intervalo pasa a ser de 0.989 a 2.31. (b) $2390 \times
12.011/44.009 = \qty{652}{PgC}$. (c) $0.45 \times 2.39 = \qty{1.08}{\celsius}$.
\end{solution}

\begin{solution}{exo:b2:anthropogenic-climate-change:5}
(a) $g = 159/398 = 0.400$; $T = (398/\num{5.670e-8})^{1/4} = \qty{289}{K}$. (b) $-0.6 \times 4
\times \num{5.670e-8} \times 288^3 = \qty{-3.25}{W.m^{-2}.K^{-1}}$; a \qty{298}{K}, $-3.60$.
(c) Menos sensible: la respuesta crece como $T^3$, así que un planeta más caliente se desprende con más facilidad de cada grado
adicional, a igualdad de las demás retroalimentaciones.
\end{solution}

\begin{solution}{exo:b2:anthropogenic-climate-change:6}
(a) $1.85 - 0.50 = \qty{1.35}{W.m^{-2}.\celsius^{-1}}$. (b) $\sqrt{0.32^2 + 0.20^2} = 0.38$.
(c) La suma se conoce mucho mejor que cada término: sus errores están fuertemente anticorrelacionados. Donde
la troposfera superior se calienta más, contiene también más vapor, de modo que una retroalimentación negativa del gradiente vertical más fuerte
va acompañada de una retroalimentación positiva del vapor de agua más fuerte.
\end{solution}

\begin{solution}{exo:b2:anthropogenic-climate-change:7}
(a) $\lambda = (0.59 - 2.20)/1.03 = \qty{-1.56}{W.m^{-2}.\celsius^{-1}}$; $3.93/1.56 =
\qty{2.52}{\celsius}$. (b) $3.93/(1.56 - 0.5) = 3.93/1.06 = \qty{3.71}{\celsius}$. (c) $3.93/1.46 =
\qty{2.69}{\celsius}$. La evaluación, con distribuciones completas, halla 3.5 y \qty{2.7}{\celsius}.
\end{solution}

\begin{solution}{exo:b2:anthropogenic-climate-change:8}
(a) $0.617/3682 = \qty{1.68e-4}{\celsius^{-1}}$. (b) $0.113 \times 0.396 = \qty{0.0447}{m}$, es decir,
\qty{44.7}{mm}, cerca de los \qty{47.5}{mm} del balance. (c) $\num{2e-4} \times 0.2 \times 700 =
\qty{0.028}{m}$, \qty{28}{mm}.
\end{solution}

\begin{solution}{exo:b2:anthropogenic-climate-change:9}
(a) $\sum xy = 0.311$, $\sum x^2 = 0.300$: $\beta = 1.04$. (b) El intervalo del \qty{95}{\percent} es
$1.04 \pm 1.96 \times 0.08$, de 0.88 a 1.20: excluye el 0, así que la respuesta está detectada, y contiene el
1, así que el cambio observado es coherente con el simulado. (c) El intervalo de 0.34 a 0.66 sigue
excluyendo el 0 (detectada), pero excluye el 1: el modelo sobrestimaría la respuesta en un factor dos
aproximadamente.
\end{solution}

\begin{solution}{exo:b2:anthropogenic-climate-change:10}
(a) $\sum g^2 = 1.20$, $\sum a^2 = 0.18$, $\sum ga = -0.46$, $\sum gy = 0.74$, $\sum ay = -0.28$.
Las ecuaciones normales son $1.20\,\beta_g - 0.46\,\beta_a = 0.74$ y $-0.46\,\beta_g + 0.18\,\beta_a
= -0.28$; su determinante es $1.20 \times 0.18 - 0.46^2 = 0.0044$, y $\beta_g = (0.74 \times
0.18 - 0.46 \times 0.28)/0.0044 = 1.00$, $\beta_a = (-1.20 \times 0.28 + 0.46 \times 0.74)/0.0044 =
1.00$. (b) Gases de efecto invernadero, $+\qty{0.80}{\celsius}$; aerosoles, $-\qty{0.30}{\celsius}$. (c) El
determinante es pequeño porque $g$ y $a$ son casi proporcionales (su correlación es $-0.99$): un
poco de ruido en $y$ puede absorberse aumentando a la vez ambos factores de escala, de modo que cada uno queda mal
determinado, aunque su efecto combinado se conozca bien.
\end{solution}

\begin{solution}{exo:b2:anthropogenic-climate-change:11}
(a) $(0.93 - 0.1)/0.45 \times 1000 = \qty{1844}{Gt}$; menos 210 y $26 \times 0.93 = 24$:
\qty{1610}{Gt}. (b) $0.93 - x = 0.45 \times (1350 + 210 + 24)/1000 = 0.713$, así que $x =
\qty{0.22}{\celsius}$: en torno al extremo superior de los 0.1 a \qty{0.2}{\celsius} evaluados. (c) El
presupuesto depende del calentamiento no debido al CO$_2$, $x$, solo a través de $-x/\text{TCRE}$: cada
\qty{0.1}{\celsius} cuesta $0.1/0.45 \times 1000 = \qty{222}{Gt}$, sea cual sea el límite.
\end{solution}

\begin{solution}{exo:b2:anthropogenic-climate-change:12}
(a) $3.93/1.81 = \qty{2.17}{\celsius}$, $3.93/0.51 = \qty{7.71}{\celsius}$, $3.93/1.16 =
\qty{3.39}{\celsius}$. (b) El lado bajo mide $3.39 - 2.17 = \qty{1.22}{\celsius}$, el lado alto
$7.71 - 3.39 = \qty{4.32}{\celsius}$: la función $1/|\lambda|$ es convexa, así que pasos iguales de
$\lambda$ hacia cero elevan la ECS cada vez más. (c) Un $\lambda$ débilmente negativo da una ECS muy
grande, y es difícil excluir pequeños cambios de $\lambda$ cerca de cero; una ECS baja requiere un $\lambda$ fuertemente
negativo, que la \omterm{def:b2:anthropogenic-climate-change:feedbacks}{respuesta de Planck} acota y el calentamiento observado contradice.
\end{solution}

\begin{solution}{pb:b2:anthropogenic-climate-change:1}
\textbf{1.} $2.16 + 0.54 + 0.21 + 0.41 = \qty{3.32}{W.m^{-2}}$.
\textbf{2.} $-0.22 - 0.84 = \qty{-1.06}{W.m^{-2}}$.
\textbf{3.} $3.32 + 0.47 + 0.05 - 1.06 - 0.20 + 0.08 + 0.06 = \qty{2.72}{W.m^{-2}}$, el total
evaluado.
\textbf{4.} $3.93\log_2(410/278) = \qty{2.20}{W.m^{-2}}$, a menos de un \qty{2}{\percent} de 2.16.
\textbf{5.} $3.93\log_2(427.35/278) = \qty{2.44}{W.m^{-2}}$.
\textbf{6.} El término aerosol--nube, de $-1.45$ a $-0.25$, con \qty{1.20}{W.m^{-2}} de anchura, frente a 1.52 para
el total: lleva la mayor parte de la incertidumbre del total porque es grande y está mal
acotado.
\textbf{7.} $0.79 \times \num{5.10e14} \times 12 \times \num{3.156e7} = \qty{1.53e23}{J}$,
\qty{153}{ZJ}.
\textbf{8.} $0.91 \times 153 = \qty{139}{ZJ}$.
\textbf{9.} $\lambda = (0.59 - 2.20)/1.03 = \qty{-1.56}{W.m^{-2}.\celsius^{-1}}$.
\textbf{10.} $3.93/1.56 = \qty{2.52}{\celsius}$.
\textbf{11.} $3.93 \times 1.03/2.20 = \qty{1.84}{\celsius}$.
\textbf{12.} $3.93/1.16 = \qty{3.39}{\celsius}$, mayor en \qty{0.87}{\celsius}.
\textbf{13.} El calentamiento ha sido hasta ahora más intenso donde las retroalimentaciones son más negativas; a medida que la
distribución evoluciona, el parámetro de retroalimentación se vuelve menos negativo, así que la sensibilidad a largo plazo
supera a la efectiva.
\textbf{14.} $0.45 \times 44.009/12.011 = \qty{1.65}{\celsius}$ por \qty{1000}{PgC}.
\textbf{15.} De 0.989 a \qty{2.31}{\celsius} por \qty{1000}{PgC}.
\textbf{16.} $2390 \times 12.011/44.009 = \qty{652}{PgC}$.
\textbf{17.} $0.45 \times 2.39 = \qty{1.08}{\celsius}$.
\textbf{18.} Dentro del redondeo, el dióxido de carbono por sí solo explica todo el calentamiento humano: el
calentamiento por los demás gases de efecto invernadero y el enfriamiento por los aerosoles casi se compensan.
\textbf{19.} $1.5 - 1.07 = \qty{0.43}{\celsius}$.
\textbf{20.} $(0.43 - 0.1)/0.45 \times 1000 = \qty{733}{Gt}$ de dióxido de carbono.
\textbf{21.} $733 - 210 = \qty{523}{Gt}$.
\textbf{22.} $26 \times 0.43 = \qty{11.2}{Gt}$; $523 - 11.2 = \qty{512}{Gt}$.
\textbf{23.} Disminuye en $0.1/0.45 \times 1000 = \qty{222}{Gt}$, hasta \qty{290}{Gt}: los
\qty{220}{Gt} en que, según la evaluación, las decisiones sobre los demás factores desplazan el presupuesto.
\textbf{24.} Calculado a partir de la TCRE, el \omterm{def:b2:anthropogenic-climate-change:budget}{presupuesto de carbono restante} para \qty{1.5}{\celsius} desde el
comienzo de 2020 es de \textbf{\qty[text-series-to-math=true]{512}{Gt}} de dióxido de carbono, a menos de un
\qty{3}{\percent} de las \qty{500}{Gt} evaluadas con una probabilidad del \qty{50}{\percent}.
\end{solution}
